%==============================================================================
%  Backpropagation:
%  Background, Theory, Chain-Rule Derivation and Worked Examples
%
%  Self-contained lecture notes for an introductory Machine Learning course
%  (B.Tech / M.Tech level).  Prerequisites: partial derivatives, the chain rule,
%  basic linear algebra.
%
%  Compile with:  pdflatex back_propagation.tex   (run twice for refs)
%  Requires a full TeX distribution (TeX Live / MiKTeX) with tikz + pgfplots.
%==============================================================================
\documentclass[11pt,a4paper]{report}

\usepackage[a4paper,margin=1in]{geometry}
\usepackage{amsmath,amssymb,amsthm,mathtools}
\usepackage{bm}
\usepackage{booktabs}
\usepackage{array}
\usepackage{enumitem}
\usepackage{caption}
\usepackage{graphicx}
\usepackage{xcolor}
\usepackage{tikz}
\usepackage{pgfplots}
\pgfplotsset{compat=1.18}
\usetikzlibrary{arrows.meta,calc,positioning}
\usepackage{listings}
\usepackage[colorlinks=true,linkcolor=blue!55!black,citecolor=blue!55!black,
            urlcolor=blue!55!black,linktocpage=true]{hyperref}

\theoremstyle{definition}
\newtheorem{definition}{Definition}[chapter]
\newtheorem{example}{Example}[chapter]
\newtheorem{exercise}{Exercise}[chapter]
\theoremstyle{remark}
\newtheorem*{remark}{Remark}

\newcommand{\R}{\mathbb{R}}
\newcommand{\vx}{\mathbf{x}}
\newcommand{\vy}{\mathbf{y}}
\newcommand{\vb}{\mathbf{b}}
\newcommand{\vz}{\mathbf{z}}
\newcommand{\va}{\mathbf{a}}
\newcommand{\vd}{\bm{\delta}}
\newcommand{\W}{\mathbf{W}}
\newcommand{\pd}[2]{\frac{\partial #1}{\partial #2}}
\DeclareMathOperator{\sig}{\sigma}

\definecolor{codegreen}{rgb}{0.0,0.5,0.0}
\definecolor{codegray}{rgb}{0.5,0.5,0.5}
\definecolor{codeblue}{rgb}{0.1,0.1,0.6}
\definecolor{backcol}{rgb}{0.97,0.97,0.95}
\definecolor{cin}{HTML}{2A9D8F}
\definecolor{chid}{HTML}{264653}
\definecolor{cout}{HTML}{E76F51}
\lstdefinestyle{python}{
  backgroundcolor=\color{backcol},language=Python,
  commentstyle=\color{codegreen},keywordstyle=\color{codeblue}\bfseries,
  numberstyle=\tiny\color{codegray},stringstyle=\color{codegreen},
  basicstyle=\ttfamily\footnotesize,breaklines=true,captionpos=b,
  numbers=left,numbersep=6pt,showstringspaces=false,frame=single,
  rulecolor=\color{codegray},tabsize=2}

%==============================================================================
\begin{document}

\begin{titlepage}
\centering
\vspace*{2cm}
{\Huge\bfseries Backpropagation\par}
\vspace{0.6cm}
{\Large Background, Theory, the Chain-Rule Derivation\\[2pt]
and Worked Examples by Hand\par}
\vspace{1.5cm}
{\large A Self-Contained Set of Lecture Notes\par}
\vspace{0.4cm}
{\large Introduction to Machine Learning\par}
\vspace{2.5cm}
\begin{tabular}{c}
\hline\\[-6pt]
\textbf{Prerequisites:} Partial Derivatives, the Chain Rule, Linear Algebra\\[4pt]
\textbf{Level:} B.Tech / M.Tech\\[4pt]
\hline
\end{tabular}
\vfill
{\large \today\par}
\end{titlepage}

\tableofcontents
\clearpage

%==============================================================================
\chapter{Background: Neural Networks in Machine Learning}
%==============================================================================

\section{From the Perceptron to the Multilayer Network}
A single \emph{artificial neuron} computes a weighted sum of its inputs, adds a
bias, and passes the result through a nonlinear \emph{activation function}:
\[
a=\varphi\!\Bigl(\sum_k w_k x_k + b\Bigr).
\]

\begin{figure}[ht]
\centering
\begin{tikzpicture}[>=Latex,every node/.style={font=\small}]
\node[circle,draw,fill=cin!20,minimum size=7mm] (x1) at (0,1.5){$x_1$};
\node[circle,draw,fill=cin!20,minimum size=7mm] (x2) at (0,0){$x_2$};
\node[circle,draw,fill=cin!20,minimum size=7mm] (x3) at (0,-1.5){$x_3$};
\node[circle,draw,thick,minimum size=11mm] (sum) at (3.3,0){$\sum$};
\node[draw,thick,minimum size=11mm] (act) at (5.8,0){$\varphi$};
\node (out) at (8.4,0){$a=\varphi(z)$};
\draw[->] (x1)--(sum) node[pos=0.55,above,sloped]{\scriptsize $w_1$};
\draw[->] (x2)--(sum) node[pos=0.5,above]{\scriptsize $w_2$};
\draw[->] (x3)--(sum) node[pos=0.55,below,sloped]{\scriptsize $w_3$};
\node[below=0.5mm of sum]{\scriptsize $+\,b$};
\draw[->] (sum)--(act) node[midway,above]{\scriptsize $z$};
\draw[->] (act)--(out);
\end{tikzpicture}
\caption{A single artificial neuron. The inputs are weighted and summed with a
bias to form the net input $z=\sum_k w_k x_k+b$, which is passed through the
activation $\varphi$ to give the output $a$.}
\label{fig:neuron}
\end{figure}
A single neuron (the \emph{perceptron}) can only separate data with a straight
line, so it cannot solve even the XOR problem. Stacking neurons into
\emph{layers}---a \textbf{multilayer perceptron} (MLP)---and using nonlinear
activations lets the network represent arbitrarily complex functions. The
question then becomes: \emph{how do we choose the weights?}

\section{Learning as Optimization}
Neural-network training is \textbf{supervised}: we have inputs with known
targets $\{(\vx^{(i)},\vy^{(i)})\}$ and a \emph{loss} $E$ measuring how wrong
the network's outputs are. Learning means adjusting every weight and bias to
make $E$ small, by \emph{gradient descent}:
\[
w \leftarrow w-\eta\,\pd{E}{w}.
\]
This needs the gradient of the loss with respect to \emph{every} weight---
possibly millions of them, buried several layers deep.
\textbf{Backpropagation} is the efficient algorithm that computes all of these
partial derivatives in a single backward sweep, using the \emph{chain rule}.

\begin{remark}[Parameters vs.\ hyperparameters]
The quantities backpropagation \emph{updates} are the \textbf{parameters}: the
weights $w$ and biases $b$. Quantities the user \emph{fixes} beforehand---the
learning rate $\eta$, the number of layers, the number of neurons, the choice
of activation---are the \textbf{hyperparameters}. Backpropagation computes the
gradients; gradient descent uses them (and $\eta$) to move the parameters.
\end{remark}

\section{Why Backpropagation Matters}
Computing each derivative independently (e.g.\ by finite differences) would cost
one forward pass \emph{per weight}---hopelessly slow. Backpropagation reuses
intermediate results so that \emph{all} gradients cost about the same as
\emph{one} extra forward pass. It is the workhorse behind essentially all modern
deep learning.

%==============================================================================
\chapter{The Multilayer Perceptron}
\label{ch:mlp}
%==============================================================================

\section{Architecture}
An MLP has an \emph{input layer}, one or more \emph{hidden layers}, and an
\emph{output layer}; every neuron in one layer connects to every neuron in the
next (``fully connected''). Figure~\ref{fig:mlp} shows a $2$--$2$--$1$ network.

\begin{figure}[ht]
\centering
\begin{tikzpicture}[
  >=Latex, shorten >=1pt, node distance=16mm and 26mm,
  neuron/.style={circle,draw,minimum size=9mm,thick},
  every node/.style={font=\small}]
% layers
\node[neuron,fill=cin!20]  (x1) at (0,1.2) {$x_1$};
\node[neuron,fill=cin!20]  (x2) at (0,-1.2){$x_2$};
\node[neuron,fill=chid!18] (h1) at (3,1.2) {$h_1$};
\node[neuron,fill=chid!18] (h2) at (3,-1.2){$h_2$};
\node[neuron,fill=cout!22] (o)  at (6,0)   {$o$};
% edges input->hidden
\foreach \i in {1,2}\foreach \j in {1,2} \draw[->] (x\i)--(h\j);
% edges hidden->output
\foreach \j in {1,2} \draw[->] (h\j)--(o);
% weight labels (selected)
\node[cin!60!black]  at ($(x1)!0.5!(h1)+(0,0.28)$) {\scriptsize $w_1$};
\node[cin!60!black]  at ($(x2)!0.5!(h1)+(-0.1,-0.28)$){\scriptsize $w_2$};
\node[cout!70!black] at ($(h1)!0.5!(o)+(0,0.30)$) {\scriptsize $w_5$};
\node[cout!70!black] at ($(h2)!0.5!(o)+(0,-0.30)$){\scriptsize $w_6$};
% layer captions
\node[above=6mm of x1] {\textbf{Input}};
\node[above=6mm of h1] {\textbf{Hidden}};
\node[above=14mm of o] {\textbf{Output}};
\end{tikzpicture}
\caption{A fully connected $2$--$2$--$1$ multilayer perceptron. Each edge carries
a weight; each hidden and output neuron also has a bias (applied inside the
neuron, not drawn).}
\label{fig:mlp}
\end{figure}

\section{Notation}
We index layers by $l=1,\dots,L$ ($l=L$ is the output). For layer $l$:
\begin{itemize}[nosep]
  \item $w^{l}_{jk}$ --- weight from the $k^{\text{th}}$ neuron in the
        $(l-1)^{\text{th}}$ layer to the $j^{\text{th}}$ neuron in the
        $l^{\text{th}}$ layer;
  \item $b^{l}_{j}$ --- bias of the $j^{\text{th}}$ neuron in the
        $l^{\text{th}}$ layer;
  \item $z^{l}_{j}=\sum_k w^{l}_{jk}\,a^{l-1}_{k}+b^{l}_{j}$ --- the
        \emph{net input} (pre-activation) of the $j^{\text{th}}$ neuron in the
        $l^{\text{th}}$ layer;
  \item $a^{l}_{j}=\varphi\!\left(z^{l}_{j}\right)$ --- the \emph{activation} of
        that neuron; with $a^{0}_{k}=x_k$ (the inputs).
\end{itemize}
In vector form, with weight matrix $\W^{l}$ and bias vector $\vb^{l}$,
\begin{equation}
\vz^{l}=\W^{l}\va^{l-1}+\vb^{l},\qquad \va^{l}=\varphi(\vz^{l}).
\label{eq:forward}
\end{equation}

\section{Forward Propagation}
\emph{Forward propagation} evaluates the network layer by layer using
\eqref{eq:forward}: start with $\va^{0}=\vx$, compute $\vz^{1},\va^{1}$, then
$\vz^{2},\va^{2}$, and so on to the output $\va^{L}$. This produces the
prediction and---crucially---stores every $z^l_j$ and $a^l_j$, which
backpropagation will reuse.

%==============================================================================
\chapter{Activation Functions}
%==============================================================================

Nonlinear activations are what give an MLP its power; their \emph{derivatives}
appear throughout backpropagation.
\begin{center}
\renewcommand{\arraystretch}{1.6}
\begin{tabular}{@{}lll@{}}
\toprule
Name & $\varphi(z)$ & $\varphi'(z)$\\
\midrule
Sigmoid & $\dfrac{1}{1+e^{-z}}$ & $\sigma(z)\bigl(1-\sigma(z)\bigr)$\\
Tanh & $\dfrac{e^{z}-e^{-z}}{e^{z}+e^{-z}}$ & $1-\tanh^2(z)$\\
ReLU & $\max(0,z)$ & $\begin{cases}1 & z>0\\0 & z<0\end{cases}$\\
Leaky ReLU & $\max(\alpha z,z)$ & $1$ if $z>0$, else $\alpha$\\
\bottomrule
\end{tabular}
\end{center}


\begin{figure}[ht]
\centering
\begin{tikzpicture}
\begin{axis}[width=0.6\textwidth,height=4.8cm,axis lines=middle,
  xlabel={$z$},ylabel={$\sigma(z)$},xmin=-6,xmax=6,ymin=-0.15,ymax=1.15,
  grid=both,major grid style={gray!20},samples=101]
\addplot[very thick,cin,domain=-6:6]{1/(1+exp(-x))};
\end{axis}
\end{tikzpicture}
\caption{\textbf{Sigmoid} $\sigma(z)=\dfrac{1}{1+e^{-z}}$ squashes any input into
the range $(0,1)$.}
\label{fig:actsigmoid}
\end{figure}

\begin{figure}[ht]
\centering
\begin{tikzpicture}
\begin{axis}[width=0.6\textwidth,height=4.8cm,axis lines=middle,
  xlabel={$z$},ylabel={$\tanh(z)$},xmin=-6,xmax=6,ymin=-1.25,ymax=1.25,
  grid=both,major grid style={gray!20},samples=101]
\addplot[very thick,chid,domain=-6:6]{tanh(x)};
\end{axis}
\end{tikzpicture}
\caption{\textbf{Tanh} is a rescaled sigmoid squashing into $(-1,1)$ and centred
at $0$.}
\label{fig:acttanh}
\end{figure}

\begin{figure}[ht]
\centering
\begin{tikzpicture}
\begin{axis}[width=0.6\textwidth,height=4.8cm,axis lines=middle,
  xlabel={$z$},ylabel={$\mathrm{ReLU}(z)$},xmin=-5,xmax=5,ymin=-0.5,ymax=5,
  grid=both,major grid style={gray!20},samples=101]
\addplot[very thick,cout,domain=-5:5]{max(0,x)};
\end{axis}
\end{tikzpicture}
\caption{\textbf{ReLU} $\max(0,z)$ passes positive inputs unchanged and clips
negatives to $0$.}
\label{fig:actrelu}
\end{figure}

\begin{definition}[Sigmoid derivative]
The sigmoid's derivative has the convenient self-referential form
\begin{equation}
\sigma(z)=\frac{1}{1+e^{-z}}\ \Longrightarrow\
\sigma'(z)=\sigma(z)\bigl(1-\sigma(z)\bigr).
\label{eq:sigder}
\end{equation}
So once $a=\sigma(z)$ is known from the forward pass, $\sigma'(z)=a(1-a)$ needs
no re-evaluation---one reason sigmoids are handy for hand computation.
\end{definition}

\begin{figure}[ht]
\centering
\begin{tikzpicture}
\begin{axis}[width=0.72\textwidth,height=5.5cm,axis lines=middle,
  xlabel={$z$},xmin=-6,xmax=6,ymin=0,ymax=1.1,
  legend pos=north west,legend cell align=left,
  grid=both,major grid style={gray!20},samples=101]
\addplot[thick,cin,domain=-6:6]{1/(1+exp(-x))};
\addlegendentry{$\sigma(z)$}
\addplot[thick,cout,dashed,domain=-6:6]{(1/(1+exp(-x)))*(1-1/(1+exp(-x)))};
\addlegendentry{$\sigma'(z)$}
\end{axis}
\end{tikzpicture}
\caption{The sigmoid and its derivative $\sigma'(z)=\sigma(z)\bigl(1-\sigma(z)\bigr)$.
The derivative peaks at only $0.25$ (at $z=0$) and is almost zero for large
$|z|$---the root of the vanishing-gradient problem.}
\label{fig:sigder}
\end{figure}

\begin{remark}[Vanishing gradients]
For large $|z|$, $\sigma'(z)\approx0$. In deep sigmoid/tanh networks the product
of many small derivatives makes early-layer gradients vanish---which is why ReLU
(derivative $1$ on the positive side) is preferred in deep models.
\end{remark}

%==============================================================================
\chapter{Loss Function and Gradient Descent}
%==============================================================================

\section{Loss Functions}
For regression, the \emph{mean squared error}; for a single example with output
vector $\va^{L}$ and target $\vy$,
\begin{equation}
E=\tfrac12\sum_{j}\bigl(a^{L}_{j}-y_j\bigr)^2 .
\label{eq:mse}
\end{equation}
The $\tfrac12$ is a convenience: it cancels the exponent $2$ on
differentiation.

For \emph{classification} the outputs are probabilities, and the natural loss is
the \textbf{cross-entropy}, which measures the disagreement between the true
label distribution and the predicted one.

\begin{definition}[Cross-entropy loss]
For \textbf{binary} classification with target $y\in\{0,1\}$ and predicted
probability $\hat y=a^{L}\in(0,1)$ (a sigmoid output),
\begin{equation}
E=-\Bigl[\,y\log\hat y+(1-y)\log(1-\hat y)\,\Bigr].
\label{eq:bce}
\end{equation}
For \textbf{multi-class} classification with $C$ classes, a one-hot target
$\vy=(y_1,\dots,y_C)$ and softmax outputs $\hat y_c=a^{L}_{c}$
(so $\sum_c\hat y_c=1$),
\begin{equation}
E=-\sum_{c=1}^{C} y_c\,\log\hat y_c .
\label{eq:cce}
\end{equation}
\end{definition}

Intuitively, cross-entropy is small when the predicted probability of the
\emph{correct} class is close to $1$ and grows without bound as that probability
approaches $0$ (a confident wrong answer is punished heavily). It also pairs
neatly with the sigmoid/softmax: the output error simplifies to the same clean
form $\delta^{L}=\hat{\vy}-\vy$, exactly as with the squared error.

\section{Gradient Descent}
Each parameter is nudged opposite its gradient:
\begin{equation}
w^{l}_{jk}\leftarrow w^{l}_{jk}-\eta\,\pd{E}{w^{l}_{jk}},\qquad
b^{l}_{j}\leftarrow b^{l}_{j}-\eta\,\pd{E}{b^{l}_{j}},
\label{eq:gd}
\end{equation}

\begin{figure}[ht]
\centering
\begin{tikzpicture}
\begin{axis}[width=0.72\textwidth,height=5.6cm,axis lines=middle,
  xlabel={weight $w$},ylabel={loss $E(w)$},
  xmin=-2.4,xmax=2.4,ymin=-0.3,ymax=4.4,xtick=\empty,ytick=\empty,samples=101]
\addplot[thick,chid,domain=-2.3:2.3]{0.8*x^2+0.2};
\addplot[only marks,mark=*,cout,mark size=2pt]
  coordinates {(2,3.4)(1.2,1.352)(0.6,0.488)(0.2,0.232)(0,0.2)};
\draw[->,thick,cout] (axis cs:2,3.4)  -- (axis cs:1.2,1.352);
\draw[->,thick,cout] (axis cs:1.2,1.352)--(axis cs:0.6,0.488);
\draw[->,thick,cout] (axis cs:0.6,0.488)--(axis cs:0.2,0.232);
\node[font=\scriptsize,anchor=west] at (axis cs:1.35,2.6)
  {$w\leftarrow w-\eta\,\dfrac{\partial E}{\partial w}$};
\node[font=\scriptsize,anchor=north] at (axis cs:0,0.2) {minimum};
\end{axis}
\end{tikzpicture}
\caption{Gradient descent on a (convex) slice of the loss. Each step moves the
weight downhill, opposite the gradient and scaled by the learning rate $\eta$,
toward the minimum.}
\label{fig:gd}
\end{figure}
where $\eta>0$ is the learning rate. Everything hinges on the partial
derivatives $\partial E/\partial w$---which backpropagation supplies.

%==============================================================================
\chapter{Backpropagation: The Chain-Rule Derivation}
\label{ch:deriv}
%==============================================================================

The loss depends on a weight only \emph{through} that neuron's net input, which
feeds its activation, which feeds the next layer, and so on. The chain rule
threads the derivative back along this path. The key trick is to define, for
each neuron, the \emph{error signal}
\begin{equation}
\boxed{\ \delta^{l}_{j}\;\coloneqq\;\pd{E}{z^{l}_{j}}\ }
\label{eq:delta}
\end{equation}
the sensitivity of the loss to that neuron's net input. Once we have every
$\delta^l_j$, the weight gradients follow immediately.

\section{Gradient of a Weight in Terms of $\delta$}
Because $z^{l}_{j}=\sum_k w^{l}_{jk}a^{l-1}_{k}+b^{l}_{j}$ depends on
$w^{l}_{jk}$ directly, the chain rule gives
\begin{equation}
\pd{E}{w^{l}_{jk}}
=\pd{E}{z^{l}_{j}}\,\pd{z^{l}_{j}}{w^{l}_{jk}}
=\delta^{l}_{j}\,a^{l-1}_{k},
\qquad
\pd{E}{b^{l}_{j}}=\delta^{l}_{j}\,\pd{z^{l}_{j}}{b^{l}_{j}}=\delta^{l}_{j}.
\label{eq:gradw}
\end{equation}
\emph{Every weight gradient is (error signal at its head) $\times$ (activation at
its tail).} It remains to compute the $\delta$'s.

\section{Error Signal at the Output Layer}
For an output neuron $j$ (layer $L$), the loss depends on $z^{L}_{j}$ only
through its own activation $a^{L}_{j}=\varphi(z^{L}_{j})$:
\begin{equation}
\delta^{L}_{j}=\pd{E}{z^{L}_{j}}
=\pd{E}{a^{L}_{j}}\,\pd{a^{L}_{j}}{z^{L}_{j}}
=\pd{E}{a^{L}_{j}}\,\varphi'(z^{L}_{j}).
\label{eq:deltaL}
\end{equation}
For the MSE loss \eqref{eq:mse}, $\partial E/\partial a^{L}_{j}=a^{L}_{j}-y_j$,
so
\begin{equation}
\boxed{\ \delta^{L}_{j}=\bigl(a^{L}_{j}-y_j\bigr)\,\varphi'(z^{L}_{j})\ }
\label{eq:deltaLmse}
\end{equation}

\section{Error Signal at a Hidden Layer (the Backward Recursion)}
A hidden neuron $j$ in layer $l$ influences the loss through \emph{all} neurons
$i$ of the next layer $l+1$ that it feeds. By the multivariate chain rule,
\begin{equation}
\delta^{l}_{j}=\pd{E}{z^{l}_{j}}
=\sum_{i}\pd{E}{z^{l+1}_{i}}\,\pd{z^{l+1}_{i}}{z^{l}_{j}}
=\sum_{i}\delta^{l+1}_{i}\,\pd{z^{l+1}_{i}}{z^{l}_{j}}.
\label{eq:hid1}
\end{equation}
Now $z^{l+1}_{i}=\sum_{j}w^{l+1}_{ij}\,\varphi(z^{l}_{j})+b^{l+1}_{i}$, so
\[
\pd{z^{l+1}_{i}}{z^{l}_{j}}=w^{l+1}_{ij}\,\varphi'(z^{l}_{j}).
\]
Substituting,
\begin{equation}
\boxed{\ \delta^{l}_{j}=\Bigl(\sum_{i}w^{l+1}_{ij}\,\delta^{l+1}_{i}\Bigr)\,\varphi'(z^{l}_{j})\ }
\label{eq:hid2}
\end{equation}
This is the heart of backpropagation: \textbf{the error signal of a neuron is the
weighted sum of the error signals it feeds, scaled by its own activation
derivative.} Errors literally \emph{propagate backward} through the same weights
used in the forward pass.

\begin{figure}[ht]
\centering
\begin{tikzpicture}[>=Latex,neuron/.style={circle,draw,minimum size=9mm,thick},
  every node/.style={font=\small}]
\node[neuron,fill=chid!18] (h1) at (0,1.3){$h_1$};
\node[neuron,fill=chid!18] (h2) at (0,-1.3){$h_2$};
\node[neuron,fill=cout!22] (o)  at (3.6,0){$o$};
\draw[->,thick,cout] (o)--(h1) node[pos=0.5,above,sloped]{\scriptsize $w_5$};
\draw[->,thick,cout] (o)--(h2) node[pos=0.5,below,sloped]{\scriptsize $w_6$};
\node[right=1.5mm of o]  {\scriptsize $\delta_o$};
\node[left=1.5mm of h1]  {\scriptsize $\delta_{h_1}$};
\node[left=1.5mm of h2]  {\scriptsize $\delta_{h_2}$};
\node[below=9mm of o,font=\small,align=center]
  {backward error flow};
\end{tikzpicture}
\caption{Backward propagation of error signals. The output error $\delta_o$ is
sent back through the \emph{same} weights $w_5,w_6$ and scaled by each hidden
neuron's activation derivative to give
$\delta_{h_1}=w_5\delta_o\,\varphi'(z_{h_1})$ and
$\delta_{h_2}=w_6\delta_o\,\varphi'(z_{h_2})$---equation \eqref{eq:hid2}.}
\label{fig:backflow}
\end{figure}

\section{Vector Form}
Collecting neurons into vectors (with $\odot$ the element-wise product),
\begin{equation}
\vd^{L}=\nabla_{\va}E\odot\varphi'(\vz^{L}),\qquad
\vd^{l}=\bigl((\W^{l+1})^{\!\top}\vd^{l+1}\bigr)\odot\varphi'(\vz^{l}),
\label{eq:vecdelta}
\end{equation}
\begin{equation}
\nabla_{\W^{l}}E=\vd^{l}\,(\va^{l-1})^{\!\top},\qquad
\nabla_{\vb^{l}}E=\vd^{l}.
\label{eq:vecgrad}
\end{equation}
The transpose $(\W^{l+1})^{\!\top}$ makes precise the phrase ``send the error
back through the weights.''

%==============================================================================
\chapter{The Backpropagation Algorithm}
%==============================================================================

\begin{center}
\fbox{\parbox{0.92\textwidth}{
\textbf{Backpropagation (one training example)}
\begin{enumerate}[nosep,leftmargin=2.4em]
  \item \textbf{Forward pass.} Set $\va^{0}=\vx$; for $l=1,\dots,L$ compute
        $\vz^{l}=\W^{l}\va^{l-1}+\vb^{l}$ and $\va^{l}=\varphi(\vz^{l})$; store
        all $\vz^{l},\va^{l}$.
  \item \textbf{Output error.} $\vd^{L}=(\va^{L}-\vy)\odot\varphi'(\vz^{L})$.
  \item \textbf{Backward pass.} For $l=L-1,\dots,1$:
        $\vd^{l}=\bigl((\W^{l+1})^{\!\top}\vd^{l+1}\bigr)\odot\varphi'(\vz^{l})$.
  \item \textbf{Gradients.} $\nabla_{\W^{l}}E=\vd^{l}(\va^{l-1})^{\!\top}$,\;
        $\nabla_{\vb^{l}}E=\vd^{l}$.
  \item \textbf{Update.} $\W^{l}\leftarrow\W^{l}-\eta\,\nabla_{\W^{l}}E$,\;
        $\vb^{l}\leftarrow\vb^{l}-\eta\,\nabla_{\vb^{l}}E$.
\end{enumerate}
}}
\end{center}
Steps 1--5 are repeated over the training data until the loss converges.

%==============================================================================
\chapter{Types and Variants}
%==============================================================================

\section{How Much Data per Update?}
\begin{center}
\begin{tabular}{@{}p{4.1cm}p{9.0cm}@{}}
\toprule
\textbf{Batch GD} & Average the gradient over the \emph{whole} training set
before each update. Stable but slow per step.\\
\textbf{Stochastic GD (SGD)} & Update after \emph{each} example. Noisy but fast;
the noise can help escape poor minima.\\
\textbf{Mini-batch GD} & Update after a small batch (e.g.\ 32--256). The
practical default---balances speed and stability.\\
\bottomrule
\end{tabular}
\end{center}

\section{Better Update Rules (Optimizers)}
Plain gradient descent \eqref{eq:gd} can be slow or unstable. Common
improvements reuse the same backprop gradients but change how they are applied:
\begin{itemize}[nosep]
  \item \textbf{Momentum} --- accumulate a velocity to smooth and accelerate.
  \item \textbf{RMSProp} --- scale each weight's step by a running average of its
        squared gradient.
  \item \textbf{Adam} --- combines momentum and RMSProp; the usual default.
\end{itemize}
Backpropagation is unchanged; only the update step differs.

%==============================================================================
\chapter{Worked Example by Hand}
\label{ch:worked}
%==============================================================================

We train the $2$--$2$--$1$ sigmoid network of Figure~\ref{fig:mlp} for one step.

\section{Setup}
\begin{center}
\begin{tabular}{@{}lll@{}}
\toprule
Inputs & $x_1=0.05,\ x_2=0.10$ & Target $t=0.01$\\
Hidden weights & $w_1=0.15,\ w_2=0.20$ (to $h_1$) & $w_3=0.25,\ w_4=0.30$ (to $h_2$)\\
Hidden biases & $b_{h_1}=b_{h_2}=0.35$ &\\
Output weights & $w_5=0.40,\ w_6=0.45$ & Output bias $b_o=0.60$\\
Activation & sigmoid ($\varphi=\sigma$) & Loss $E=\tfrac12(o-t)^2$,\ \ $\eta=0.5$\\
\bottomrule
\end{tabular}
\end{center}
Here $w_1,w_2$ feed $h_1$; $w_3,w_4$ feed $h_2$; $w_5,w_6$ feed $o$.

\section{Step 1 --- Forward Pass}
\begin{align*}
z_{h_1}&=w_1x_1+w_2x_2+b_{h_1}=0.15(0.05)+0.20(0.10)+0.35=0.3775,\\
h_1&=\sigma(0.3775)=0.593270,\\
z_{h_2}&=w_3x_1+w_4x_2+b_{h_2}=0.25(0.05)+0.30(0.10)+0.35=0.3925,\\
h_2&=\sigma(0.3925)=0.596884,\\
z_{o}&=w_5h_1+w_6h_2+b_o=0.40(0.593270)+0.45(0.596884)+0.60=1.105906,\\
o&=\sigma(1.105906)=0.751365,\qquad
E=\tfrac12(0.751365-0.01)^2=0.274811.
\end{align*}

\section{Step 2 --- Output Error and its Gradients}
Using \eqref{eq:deltaLmse} with $\sigma'(z_o)=o(1-o)$:
\[
\delta_o=(o-t)\,o(1-o)=(0.751365-0.01)(0.751365)(0.248635)=0.138499.
\]
Then by \eqref{eq:gradw}, $\partial E/\partial w_5=\delta_o h_1$, etc.:
\[
\pd{E}{w_5}=0.138499(0.593270)=0.082167,\quad
\pd{E}{w_6}=0.138499(0.596884)=0.082668,\quad
\pd{E}{b_o}=\delta_o=0.138499.
\]

\section{Step 3 --- Hidden Errors and their Gradients}
By the backward recursion \eqref{eq:hid2}, with only one output neuron,
$\delta_{h}=\delta_o\,w_{(h\to o)}\,h(1-h)$:
\begin{align*}
\delta_{h_1}&=\delta_o\,w_5\,h_1(1-h_1)=0.138499(0.40)(0.593270)(0.406730)=0.0133679,\\
\delta_{h_2}&=\delta_o\,w_6\,h_2(1-h_2)=0.138499(0.45)(0.596884)(0.403116)=0.0149961.
\end{align*}
Weight gradients $\partial E/\partial w=\delta_h\,x$ (input at the tail):
\[
\begin{aligned}
\pd{E}{w_1}&=\delta_{h_1}x_1=0.0006684, & \pd{E}{w_2}&=\delta_{h_1}x_2=0.0013368, & \pd{E}{b_{h_1}}&=\delta_{h_1}=0.013368,\\
\pd{E}{w_3}&=\delta_{h_2}x_1=0.0007498, & \pd{E}{w_4}&=\delta_{h_2}x_2=0.0014996, & \pd{E}{b_{h_2}}&=\delta_{h_2}=0.014996.
\end{aligned}
\]

\section{Step 4 --- Update the Parameters ($\eta=0.5$)}
Apply $w\leftarrow w-\eta\,\partial E/\partial w$:
\begin{center}
\renewcommand{\arraystretch}{1.2}
\begin{tabular}{@{}cccc|cccc@{}}
\toprule
Param & old & grad & new & Param & old & grad & new\\
\midrule
$w_1$ & $0.15$ & $0.000668$ & $0.149666$ & $w_5$ & $0.40$ & $0.082167$ & $0.358916$\\
$w_2$ & $0.20$ & $0.001337$ & $0.199332$ & $w_6$ & $0.45$ & $0.082668$ & $0.408666$\\
$w_3$ & $0.25$ & $0.000750$ & $0.249625$ & $b_{h_1}$ & $0.35$ & $0.013368$ & $0.343316$\\
$w_4$ & $0.30$ & $0.001500$ & $0.299250$ & $b_{h_2}$ & $0.35$ & $0.014996$ & $0.342502$\\
      &        &            &            & $b_o$ & $0.60$ & $0.138499$ & $0.530751$\\
\bottomrule
\end{tabular}
\end{center}

\section{Step 5 --- Verify the Loss Decreased}
Running a forward pass with the updated weights gives a new output
$o'=0.728352$ and
\[
E'=\tfrac12(0.728352-0.01)^2=0.258015\ <\ 0.274811=E.
\]
One backpropagation step reduced the error, exactly as gradient descent
promises. Repeating the five steps drives $o$ toward the target $t=0.01$.

\begin{remark}[Why the output weights changed far more than the input weights]
$\partial E/\partial w_5\approx0.082$ but $\partial E/\partial w_1\approx0.0007$.
The hidden gradients carry the extra factors $w_5\,h_1(1-h_1)$ and the tiny
input $x_1=0.05$, shrinking them---an early glimpse of how gradients get smaller
as we move back through the network (the vanishing-gradient effect).
\end{remark}

%==============================================================================
\chapter{Practical Notes, Pros \& Cons, and Code}
%==============================================================================

\section{Practical Notes}
\begin{itemize}[nosep]
  \item \textbf{Weight initialization} matters: start small and random (e.g.\
        Xavier/He) to break symmetry and avoid saturating activations.
  \item \textbf{Learning rate} $\eta$ is the key hyperparameter: too large
        diverges, too small crawls.
  \item \textbf{Feature scaling} of inputs speeds convergence.
  \item \textbf{Vanishing/exploding gradients} plague deep networks; ReLU,
        normalization and residual connections help.
\end{itemize}

\section{Advantages and Disadvantages}
\begin{center}
\begin{tabular}{@{}p{6.4cm}p{6.4cm}@{}}
\toprule
\textbf{Advantages} & \textbf{Disadvantages}\\
\midrule
Computes all gradients in one backward pass ($\mathcal{O}$ of one forward pass).
& Needs differentiable activations and loss.\\
Exact gradients (not approximations). & Converges only to a local minimum; $E$ is
non-convex.\\
Scales to millions of parameters. & Sensitive to $\eta$, initialization,
vanishing gradients.\\
\bottomrule
\end{tabular}
\end{center}

\section{From-Scratch Implementation (NumPy)}
\begin{lstlisting}[style=python,caption={The worked example in NumPy.}]
import numpy as np
sig  = lambda z: 1/(1+np.exp(-z))
dsig = lambda a: a*(1-a)              # sigma'(z) written in terms of a=sigma(z)

x = np.array([[0.05],[0.10]]); t = np.array([[0.01]]); eta = 0.5
W1 = np.array([[0.15,0.20],[0.25,0.30]]); b1 = np.array([[0.35],[0.35]])
W2 = np.array([[0.40,0.45]]);             b2 = np.array([[0.60]])

# forward
a1 = sig(W1@x + b1)                  # hidden activations
a2 = sig(W2@a1 + b2)                 # output
E  = 0.5*np.sum((a2-t)**2)

# backward
d2 = (a2-t)*dsig(a2)                 # output error signal
d1 = (W2.T@d2)*dsig(a1)              # hidden error signal
gW2, gb2 = d2@a1.T, d2               # gradients
gW1, gb1 = d1@x.T,  d1

# update
W2 -= eta*gW2; b2 -= eta*gb2
W1 -= eta*gW1; b1 -= eta*gb1
print("loss:", E)                    # 0.274811...
\end{lstlisting}
This reproduces the by-hand numbers: $\delta_o=0.1385$, updated
$w_5=0.3589$, and so on.

%==============================================================================
\chapter{Exercises}
%==============================================================================

\section{Solved Problems}
\begin{example}[Sigmoid derivative]
If $a=\sigma(z)=0.75$, find $\sigma'(z)$.\\
\textbf{Solution.} $\sigma'(z)=a(1-a)=0.75(0.25)=0.1875.$
\end{example}
\begin{example}[Output error signal]
For an output $o=0.8$, target $t=1$, MSE loss, sigmoid: find $\delta_o$.\\
\textbf{Solution.} $\delta_o=(o-t)o(1-o)=(-0.2)(0.8)(0.2)=-0.032.$
\end{example}
\begin{example}[Weight gradient]
A neuron with error signal $\delta=0.05$ receives activation $a=0.6$ along a
weight $w$. Find $\partial E/\partial w$ and the update for $\eta=0.1$.\\
\textbf{Solution.} $\partial E/\partial w=\delta a=0.03$; $w\leftarrow w-0.1(0.03)=w-0.003.$
\end{example}

\section{Unsolved Problems}
\begin{exercise}[Second iteration]
Using the updated weights from Chapter~\ref{ch:worked}, perform a second forward
and backward pass and confirm the loss decreases again.
\end{exercise}
\begin{exercise}[Two outputs]
Extend the worked example to two output neurons with targets $(0.01,0.99)$ and
derive $\delta_{o_1},\delta_{o_2}$.
\end{exercise}
\begin{exercise}[ReLU hidden layer]
Redo Step~3 assuming the hidden neurons use ReLU instead of sigmoid. How does
$\delta_h$ change when $z_h>0$?
\end{exercise}
\begin{exercise}[Chain rule]
Starting from $E=\tfrac12(o-t)^2$ and $o=\sigma(z_o)$, derive
$\partial E/\partial w_5=\delta_o h_1$ from first principles.
\end{exercise}
\begin{exercise}[Vanishing gradient]
Show that if every $\sigma'(z)\le 0.25$ (true for sigmoid) and all $|w|\le 1$,
then $|\delta^{l}|$ shrinks by at least a factor $4$ per layer going backward.
\end{exercise}

%==============================================================================
\appendix
\chapter{Notation Summary}
%==============================================================================
\begin{center}
\begin{tabular}{@{}ll@{}}
\toprule
Symbol & Meaning\\
\midrule
$L$ & number of layers (output is layer $L$)\\
$w^{l}_{jk}$ & weight: $k^{\text{th}}$ neuron (layer $l-1$) to $j^{\text{th}}$ neuron (layer $l$)\\
$b^{l}_{j}$ & bias of the $j^{\text{th}}$ neuron in the $l^{\text{th}}$ layer\\
$z^{l}_{j}$ & net input of the $j^{\text{th}}$ neuron in the $l^{\text{th}}$ layer\\
$a^{l}_{j}$ & activation $\varphi(z^{l}_{j})$; $a^{0}=x$ (inputs)\\
$\varphi,\varphi'$ & activation function and its derivative\\
$E$ & loss (error)\\
$\delta^{l}_{j}=\partial E/\partial z^{l}_{j}$ & error signal of the $j^{\text{th}}$ neuron in the $l^{\text{th}}$ layer\\
$\eta$ & learning rate (a hyperparameter)\\
$\odot$ & element-wise (Hadamard) product\\
\bottomrule
\end{tabular}
\end{center}

\end{document}
